\documentclass[10pt,conference]{IEEEtran}

\usepackage{amsmath,amssymb}
\usepackage{graphicx}
\usepackage{booktabs}
\usepackage{array}
\usepackage{textcomp}
\usepackage{xcolor}
\usepackage{listings}
\usepackage[hidelinks]{hyperref}

\lstset{
  basicstyle=\ttfamily\scriptsize,
  breaklines=true,
  columns=fullflexible,
  frame=single,
  language=Python,
  backgroundcolor=\color{black!3},
  showstringspaces=false
}

\newcommand{\kalbee}{\textsc{kalbee}}

\begin{document}
\renewcommand{\arraystretch}{0.85}

\title{\kalbee: A Modular, Batteries-Included Python Library\\ for Kalman Filtering and State Estimation}

\author{
\IEEEauthorblockN{Le Duc Minh}
\IEEEauthorblockA{Independent Developer\\
Email: minh.leduc.0210@gmail.com\\
Repository: \url{https://github.com/MinLee0210/kalbee}}
}

\maketitle

\begin{abstract}
State estimation via the Kalman filter and its many nonlinear and robust
variants remains one of the most widely used tools in robotics, navigation,
finance, and signal processing, yet the Python ecosystem for it is
fragmented: existing libraries are either narrowly scoped, largely
unmaintained, or so heavyweight that they impose a steep learning curve for
routine tracking problems. This paper introduces \kalbee{}, an open-source
(Apache-2.0) Python library that unifies eighteen filter variants, four
smoothers, multi-target tracking, parameter-learning routines, and
applied sensor-fusion recipes behind a single, consistent
\texttt{BaseFilter} interface. We survey the closest prior art
(FilterPy, pykalman, simdkalman, Stone Soup, and JAX-based differentiable
filters), describe \kalbee{}'s architecture and numerical-stability
guarantees, and report a reproducible, head-to-head benchmark against three
of these libraries on an identical tracking task. The results show
\kalbee{} reproduces reference implementations' accuracy exactly while its
batched \texttt{VectorizedKalmanFilter} outperforms simdkalman's headline
vectorization pitch by approximately $4\times$ on 1{,}000 concurrent
series, at the cost of higher per-call overhead than a minimal
single-filter implementation. We also describe two applied recipes
--- a quaternion attitude EKF with an exact process Jacobian, and a
loosely-coupled GPS+IMU fusion filter --- and the developer-facing tooling
(a scikit-learn-compatible estimator and a command-line interface) built
to lower the barrier to adoption.
\end{abstract}

\begin{IEEEkeywords}
Kalman filter, extended Kalman filter, unscented Kalman filter, state
estimation, sensor fusion, multi-target tracking, open-source software,
Python
\end{IEEEkeywords}

\section{Introduction}

The Kalman filter~\cite{kalman1960} and its descendants --- the extended
Kalman filter (EKF), the unscented Kalman filter
(UKF)~\cite{julier2004unscented}, the cubature Kalman filter
(CKF)~\cite{arasaratnam2009cubature}, and a long tail of robust and
adaptive variants --- are the default tool for recursive state estimation
under uncertainty. Applications span inertial navigation, GPS/IMU fusion,
multi-object tracking, computational finance, and, increasingly, as a
differentiable component inside learned perception pipelines. Despite this
ubiquity, a practitioner reaching for Python today faces a fragmented
choice:

\begin{itemize}
\item \textbf{FilterPy}~\cite{filterpy} is the most widely known option,
  built as the companion code for the pedagogical text \emph{Kalman and
  Bayesian Filters in Python}. Its implementations favor a 1:1 mapping to
  textbook equations over performance or breadth, and the project has seen
  little maintenance activity in recent years.
\item \textbf{pykalman}~\cite{pykalman} offers a minimal
  KF/UKF/EM interface with a scikit-learn-flavored API, but has not added a
  new filter type in years and does not track modern NumPy releases
  smoothly.
\item \textbf{simdkalman}~\cite{simdkalman} solves exactly one problem ---
  vectorizing a linear KF across many independent series --- very well, but
  offers no nonlinear filters, smoothers, or tracking utilities.
\item \textbf{Stone Soup}~\cite{stonesoup}, maintained by the UK Defence
  Science and Technology Laboratory (Dstl), is the most comprehensive
  option: a full research framework for tracking, sensor management, and
  data association. Its generality comes at the cost of a much larger
  conceptual surface (trackers, hypothesisers, readers, feeders, and
  component registries) that is arguably disproportionate for a developer
  who just needs to fuse a GPS fix with an accelerometer reading.
\item Differentiable, GPU-oriented implementations such as
  \textbf{kalman-jax}~\cite{kalmanjax} target the machine-learning research
  community and require adopting the JAX ecosystem wholesale.
\end{itemize}

No existing library combines (i) broad algorithmic coverage, (ii) a single
consistent interface across all filter types, (iii) numerically robust
default implementations, (iv) applied, ready-to-use sensor-fusion recipes,
and (v) developer-facing tooling that lowers the barrier from ``I found
this on GitHub'' to ``I have a working tracker.'' \kalbee{} is designed to
fill that gap.

This paper makes the following contributions:
\begin{enumerate}
\item A survey and feature comparison of the current Python Kalman-filtering
  ecosystem (Section~\ref{sec:related}).
\item A description of \kalbee{}'s architecture, its common filter
  interface, and the numerical-stability techniques applied uniformly
  across all eighteen filter implementations (Section~\ref{sec:design}).
\item Two applied sensor-fusion recipes --- a quaternion attitude EKF with
  an exact (non-linearized) process Jacobian, and a loosely-coupled GPS+IMU
  fusion filter --- packaged as tested, documented building blocks rather
  than one-off scripts (Section~\ref{sec:cookbook}).
\item A reproducible, honest benchmark against FilterPy, pykalman, and
  simdkalman on an identical task, including a case where \kalbee{}
  underperforms a minimal alternative (Section~\ref{sec:eval}).
\end{enumerate}

\section{Related Work}
\label{sec:related}

Table~\ref{tab:comparison} summarizes the feature coverage of the five
systems discussed above alongside \kalbee{}. We expand on each below.

\textbf{FilterPy}~\cite{filterpy} implements the KF, EKF, UKF, a
Rauch--Tung--Striebel (RTS) smoother~\cite{rauch1965maximum}, a particle
filter, an ensemble Kalman filter, and $g$-$h$/$\alpha$-$\beta$ filters. Its
code intentionally prioritizes matching textbook notation over
vectorization or defensive numerics, which makes it an excellent teaching
tool but a less suitable foundation for production tracking systems that
must run for long, unattended horizons without covariance collapse or
divergence.

\textbf{pykalman}~\cite{pykalman} exposes a compact
\texttt{KalmanFilter} class with \texttt{filter}/\texttt{smooth}/\texttt{em}
methods in an interface reminiscent of scikit-learn estimators. It covers
only the linear KF, the UKF, and EM-based parameter learning; it has no
concept of multi-target tracking, robust filtering, or square-root/
information-form numerics.

\textbf{simdkalman}~\cite{simdkalman} vectorizes the KF predict/update
recursion across the leading axis of a NumPy array, so that $N$ independent
linear KF instances execute as a handful of batched matrix operations
instead of $N$ Python-level loop iterations. This is a genuinely useful,
narrow contribution — the project explicitly does not attempt to cover
nonlinear filtering, smoothing, or tracking.

\textbf{Stone Soup}~\cite{stonesoup} is architected around composable
components (predictors, updaters, hypothesisers, data associators, feeders,
metrics) intended for tracking-algorithm \emph{research}, including
multi-target tracking with the Probability Hypothesis Density (PHD) filter,
Gromov flow particle filters, and sensor-management strategies. It is, by
design, optimized for flexibility and experimental comparison rather than
for the ergonomics of a developer who wants to \texttt{pip install} a
tracker and start filtering within a few lines of code.

\textbf{kalman-jax}~\cite{kalmanjax} and related differentiable filters
target Markov Gaussian-process inference with JIT compilation and
autodiff, letting a user backpropagate through a filtering pipeline (e.g.,
to learn $Q$/$R$ via gradient descent, or run on GPU/TPU). This class of
tool assumes familiarity with JAX and is not intended as a general-purpose
filtering library.

\begin{table*}[t]
\centering
\caption{Feature comparison of Python Kalman-filtering libraries (as of September 2026).}
\label{tab:comparison}
\begin{tabular}{lcccccc}
\toprule
\textbf{Feature} & \textbf{\kalbee{}} & \textbf{FilterPy} & \textbf{pykalman} & \textbf{simdkalman} & \textbf{Stone Soup} & \textbf{kalman-jax} \\
\midrule
Filter implementations                  & \textbf{18}         & $\sim$10            & 2                & 1                & $\sim$6              & 1 (GP-Kalman) \\
Multi-object tracking (SORT/JPDA/PMBM)  & \checkmark           & --                & --               & --               & \checkmark (heavier)  & -- \\
Smoothers                               & RTS/EKF/UKF/fixed-lag & RTS only          & RTS only         & --               & \checkmark            & \checkmark \\
Parameter learning (EM, auto-tune)      & \checkmark           & --                & EM only          & --               & partial              & gradient-based \\
Neural hybrid filter                    & \checkmark (KalmanNet) & --              & --               & --               & --                   & -- \\
Vectorized / batched filtering          & \checkmark           & --                & --               & \checkmark        & --                   & \checkmark (JAX \texttt{vmap}) \\
DataFrame (pandas/Polars) integration   & \checkmark           & --                & --               & --               & --                   & -- \\
scikit-learn \texttt{fit}/\texttt{transform} API & \checkmark  & --                & partial          & --               & --                   & -- \\
Command-line interface                  & \checkmark           & --                & --               & --               & --                   & -- \\
GPU / autodiff backend                  & --                   & --                & --               & --               & --                   & \checkmark \\
Actively maintained (2026)              & \checkmark           & mostly dormant    & mostly dormant   & mostly dormant   & \checkmark            & sporadic \\
\bottomrule
\end{tabular}
\end{table*}

The gap this table makes visible is not any single missing filter, but the
absence of a library that is simultaneously broad, well-tested, and
approachable. Stone Soup covers similar breadth but at a much higher
conceptual cost; the narrower libraries (FilterPy, pykalman, simdkalman)
each solve one slice of the problem well but none generalize.

\section{Design and Architecture}
\label{sec:design}

\subsection{Common Filter Interface}

Every filter in \kalbee{} subclasses an abstract \texttt{BaseFilter} that
defines a uniform contract:

\begin{lstlisting}
class BaseFilter(ABC):
    def predict(self, dt=1.0, **kwargs): ...
    def update(self, measurement, **kwargs): ...
    def measure(self, state=None): ...
    def predict_only(self, dt=1.0, **kwargs): ...
    def reset(self, state=None, covariance=None): ...
    def filter_sequence(self, measurements,
                         dt=1.0, missing=None): ...
    def save_state(self, filepath): ...
    def load_state(self, filepath): ...
\end{lstlisting}

Because every filter --- linear, nonlinear, particle-based, or
information-form --- exposes the same \texttt{predict}/\texttt{update}
pair, application code can swap a \texttt{KalmanFilter} for an
\texttt{AdaptiveKalmanFilter} or a \texttt{SquareRootKalmanFilter} without
modification. A factory class, \texttt{AutoFilter}, additionally allows
filter selection by string name (e.g., \texttt{mode="ekf"}), which the
CLI and the scikit-learn wrapper both build on (Section~\ref{sec:dx}).

\subsection{Numerical Stability}

Long-running or high-noise filtering is prone to covariance
non-positive-definiteness from floating-point round-off. \kalbee{} applies
three defenses uniformly:
\begin{itemize}
\item \textbf{Joseph-form covariance updates} in the KF, EKF, and
  H-infinity filters, which are guaranteed positive semi-definite for any
  gain $K$, rather than the algebraically-equivalent but numerically
  fragile $P \leftarrow (I - KH)P$ form.
\item \textbf{Symmetry enforcement}, $P \leftarrow (P + P^{\top})/2$, after
  every covariance update.
\item \textbf{Regularized/pseudo-inverse fallbacks} (\texttt{safe\_inv},
  \texttt{safe\_cholesky} in \texttt{kalbee.modules.utils.linalg}) for
  near-singular innovation covariances, plus a dedicated
  \texttt{SquareRootKalmanFilter} and \texttt{SquareRootUKF} that propagate
  the Cholesky factor of $P$ directly for applications where
  positive-definiteness must never be violated.
\end{itemize}

\subsection{Extensibility: The Sigma-Point Strategy Pattern}

Rather than hard-coding one sigma-point rule inside the UKF, \kalbee{}
factors sigma-point generation into a pluggable
\texttt{SigmaPoints} strategy (Simplex, Merwe-scaled, and Julier
variants), so that \texttt{SigmaPointUKF(..., sigma\_points=} \texttt{MerweScaledSigmaPoints(n))}
and its Simplex/Julier counterparts share one filter implementation. Adding
a new filter follows the same pattern used throughout the codebase:
subclass \texttt{BaseFilter}, implement \texttt{predict}/\texttt{update},
and register the new class with \texttt{AutoFilter}.

\subsection{Module Layout}

The package is organized by concern: \texttt{modules/filters} (18 filter
classes plus the procedural functional API), \texttt{modules/smoothers}
(RTS, extended RTS, unscented RTS, fixed-lag), \texttt{modules/fusion}
(covariance intersection, track-to-track fusion, an asynchronous
out-of-sequence-measurement buffer), \texttt{modules/learning} (offline/
online EM, NIS-based auto-tuning, a PyTorch \texttt{KalmanNet} hybrid),
\texttt{modules/integration} (pandas, Polars, and scikit-learn adapters),
\texttt{tracking} (multi-object tracking and data association), and
\texttt{models} (ready-made motion/measurement models and the sensor-fusion
cookbook described in Section~\ref{sec:cookbook}).

\section{Algorithmic Coverage}
\label{sec:algo}

Table~\ref{tab:filters} enumerates \kalbee{}'s eighteen filter
implementations. Beyond single-target filtering, \kalbee{} provides:

\textbf{Smoothing.} A standard RTS smoother for linear systems, an
extended RTS smoother that linearizes at each backward step, an unscented
RTS smoother, and a fixed-lag smoother for bounded-delay online smoothing.

\textbf{Multi-target tracking.} A SORT-style~\cite{bewley2016sort}
\texttt{MultiObjectTracker} performs Hungarian-algorithm association with
configurable gating, track initiation (\texttt{n\_init}), and track deletion
(\texttt{max\_age}); \texttt{JPDAAssociation} implements joint probabilistic
data association for cluttered scenes; and \texttt{PMBMTracker} implements
a Poisson multi-Bernoulli mixture filter for principled track existence and
birth/death modeling under heavy clutter.

\textbf{Parameter learning.} Offline expectation-maximization
(\texttt{em\_kalman}) fits $Q$/$R$ from a batch of measurements; an online
EM variant adapts them from a stream without retaining full history; and
NIS-based auto-tuning (\texttt{tune\_kalman\_filter}, \texttt{quick\_tune})
adjusts $Q$/$R$ so that the mean normalized innovation squared matches the
measurement dimension, the standard consistency criterion for a
well-tuned filter. A PyTorch \texttt{KalmanNet} module additionally allows
the gain to be learned end-to-end when the process model is only
approximately known.

\begin{table}[t]
\centering
\caption{Filter implementations in \kalbee{}.}
\label{tab:filters}
\small
\begin{tabular}{p{0.32\linewidth}p{0.58\linewidth}}
\toprule
\textbf{Filter} & \textbf{Notes} \\
\midrule
KF & Standard linear Kalman filter, Joseph form \\
EKF & First-order linearization via user-supplied Jacobians \\
UKF / SigmaPointUKF & Unscented transform; pluggable sigma-point strategy \\
CKF & Third-order spherical-radial cubature rule \\
Particle Filter & Sequential Monte Carlo, non-Gaussian posteriors \\
RBPF & Rao-Blackwellized PF, marginalizes linear substates \\
Ensemble KF & Ensemble-based, high-dimensional systems \\
Information Filter & Dual (inverse-covariance) form; multi-sensor fusion \\
Alpha-Beta-Gamma & Constant-gain $O(1)$ tracker \\
Adaptive KF & Online $Q$/$R$ estimation (IAE) \\
Square-Root KF / UKF & Propagates the Cholesky factor of $P$ \\
Cholesky KF & Cholesky-factor-based covariance update \\
Vectorized KF & Batches $N$ independent filters via NumPy broadcasting \\
Fading Memory KF & Discounted covariance inflation, maneuvering targets \\
H-Infinity & Robust, worst-case-bounded estimation \\
IMM & Interacting Multiple Model (mode-switching) \\
Invariant EKF & Error dynamics on SO(3)/SE(3) Lie groups \\
VBAKF & Variational-Bayes adaptive KF \\
Async KF & \texttt{async}/\texttt{await} interface for streaming sensors \\
\bottomrule
\end{tabular}
\end{table}

\section{Applied Sensor-Fusion Recipes}
\label{sec:cookbook}

A recurring pattern observed across dozens of unmaintained public
repositories is a one-off script that fuses a gyroscope and accelerometer,
or a GPS receiver and an accelerometer, hand-derived from scratch each
time. \kalbee{} packages two such recipes as tested library functions.

\subsection{Quaternion Attitude EKF}

Orientation is estimated as a unit quaternion
$q = [w,x,y,z]^{\top}$ (Hamilton convention, body-to-world). Given a
body-frame angular-velocity reading $\boldsymbol\omega$ over a step
$\Delta t$, the exact integral of constant angular velocity is the delta
quaternion
\begin{equation}
\Delta q = \left[\cos\tfrac{\theta}{2},\ \hat{\boldsymbol\omega}\sin\tfrac{\theta}{2}\right]^{\top},
\qquad \theta = \lVert\boldsymbol\omega\rVert\,\Delta t,
\end{equation}
and the process model is the Hamilton product $q_{k+1} = q_k \otimes
\Delta q$. Because $\Delta q$ does not depend on $q_k$, this update is
\emph{linear} in the state, and its EKF process Jacobian is the exact
right-multiplication matrix $R(\Delta q)$ of $\Delta q$ --- not a
first-order approximation, which distinguishes it from most EKF process
Jacobians:
\begin{equation}
F_k = \frac{\partial q_{k+1}}{\partial q_k} = R(\Delta q).
\end{equation}
The measurement model predicts the (unit) accelerometer reading as gravity
rotated into the body frame, $h(q) = R(q)^{\top}\,[0,0,1]^{\top}$, with an
analytic Jacobian cross-checked in the test suite against a
finite-difference Jacobian (Section~\ref{sec:dx}).

Two properties of this filter are worth stating explicitly because they
are easy to get wrong and are rarely documented in the many ad hoc
implementations available online. First, gravity alone cannot observe
rotation about the vertical (yaw) axis --- any two orientations differing
only in yaw predict the same accelerometer reading, so gyroscope noise on
that axis random-walks uncorrected unless a magnetometer or GPS-course
update is added. Second, because this is a \emph{naive} quaternion-state
EKF (the raw four-vector is the filter state) rather than an error-state/
multiplicative EKF, its linearization is only locally valid; supplying the
raw accelerometer noise variance as $R$ allows single updates to overshoot
the unit-quaternion manifold before renormalization, which degrades
convergence markedly. In our test scenario (3-axis rotation at
$\sim$0.05--0.15\,rad/s, 5\,s at 100\,Hz), using the literal sensor noise
variance ($R = 0.02^2 I$) produced a final orientation error that grew
past $20^{\circ}$, while inflating $R$ to $0.05\,I$ --- a design choice, not
a change to the physical noise --- reduced the same scenario's error to
under $1^{\circ}$ across multiple random seeds. We report this not as a
flaw to hide but as exactly the kind of tuning knowledge a reusable
library should encode once, in documentation and tests, rather than
leaving each downstream user to rediscover independently.

\subsection{Loosely-Coupled GPS + IMU Fusion}

Rather than implementing a dedicated inertial-navigation mechanization,
\kalbee{} reuses the existing control-input term of the standard KF predict
step, $x \leftarrow Fx + Bu$. A helper,
\texttt{imu\_velocity\_control(dt, n\_dims)}, builds
\begin{equation}
B = \begin{bmatrix} \tfrac{1}{2}\Delta t^2 \\ \Delta t \end{bmatrix}
\ \text{(per axis, block-diagonal across axes)},
\end{equation}
so that feeding a gravity-compensated, world-frame accelerometer reading
as $u$ at every IMU tick advances position and velocity exactly as a
constant-acceleration integration step would, while a GPS position fix
arriving at a much lower rate triggers a standard measurement update
through the existing \texttt{position\_measurement\_model} builder. This
reduces GPS+IMU fusion to $\sim$10 lines of application code with no
custom filter subclass required.

\section{Developer Experience and Tooling}
\label{sec:dx}

Beyond algorithmic coverage, \kalbee{} includes several features aimed
specifically at lowering the time-to-first-success for a new user.

\textbf{\texttt{numerical\_jacobian}.} Deriving an EKF's Jacobians by hand
is a common source of setup bugs. \kalbee{} provides a central
finite-difference Jacobian utility that turns any
\texttt{transition\_function}/\texttt{measurement\_function} into a working
\texttt{F}/\texttt{H} without manual differentiation, at the cost of one
extra function evaluation per state dimension per step.

\textbf{scikit-learn integration.} \texttt{KalmanEstimator} wraps any
\kalbee{} filter behind \texttt{fit}/\texttt{transform}/\texttt{predict},
so it composes with \texttt{sklearn.pipeline.Pipeline} and
\texttt{GridSearchCV}. Each \texttt{transform} call resets the underlying
filter to its initial state before running, so repeated or parallel calls
do not leak state --- a subtlety that a naive wrapper around a stateful
filter object would otherwise get wrong.

\textbf{Command-line interface.} \texttt{kalbee demo}, \texttt{kalbee
bench}, and \texttt{kalbee new} let a user see a filter tracking a signal,
compare filter speed/accuracy, or scaffold a starter script without
writing any code. The \texttt{--live} flag renders an animated,
Unicode-block-sparkline view of the true signal, the noisy measurement,
and the filtered estimate directly in the terminal, updated every
predict/update step via the \texttt{rich} library.

\textbf{Typed public API.} The package ships a \texttt{py.typed} marker
(PEP 561) and a \texttt{mypy} configuration, giving IDE autocompletion and
static type checking for the public API surface.

\section{Experimental Evaluation}
\label{sec:eval}

\subsection{Methodology}

We construct a single shared task --- tracking a noisy sine-wave position
signal with a constant-velocity model ($\Delta t = 0.05$\,s, 200 steps,
measurement noise $\sigma = 0.3$, process-noise scale $0.1$) --- and
implement it identically (same $F$, $Q$, $H$, $R$, and measurement
sequence) against \kalbee{}, FilterPy, pykalman, and simdkalman. The
comparison script is published alongside the library so the results below
are independently reproducible.\footnote{\texttt{scripts/compare\_benchmarks.py} in the \kalbee{} repository.}
All timings are wall-clock means over multiple repetitions on the same
machine in the same process.

\subsection{Correctness}

All three linear-KF implementations (\kalbee{}, FilterPy, pykalman)
produce \emph{bit-for-bit identical} position RMSE (0.6380) on the shared
task. This is a useful cross-check: it confirms \kalbee{}'s
\texttt{KalmanFilter} implements the standard predict/update equations
correctly rather than merely ``close enough.''

\subsection{Single-Series Speed}

Table~\ref{tab:bench-single} reports per-run wall-clock time for a single
filter stepping through the 200-sample series once.

\begin{table}[t]
\centering
\caption{Single-series timing on the shared tracking task.}
\label{tab:bench-single}
\begin{tabular}{lrr}
\toprule
\textbf{Library} & \textbf{Time / run (ms)} & \textbf{Position RMSE} \\
\midrule
FilterPy   & 6.9  & 0.6380 \\
\kalbee{}  & 12.5 & 0.6380 \\
simdkalman & 25.5 & 0.6745 \\
pykalman   & 34.1 & 0.6380 \\
\bottomrule
\end{tabular}
\end{table}

FilterPy's minimal, single-purpose implementation has lower per-call
overhead than \kalbee{} for one bare \texttt{predict}/\texttt{update} loop
--- unsurprising, since FilterPy does strictly less (no batching support,
no shared abstract base class dispatch, seventeen fewer filter types to
keep interoperable). \kalbee{} is nonetheless faster than both pykalman and
simdkalman's single-series code path in this setting.

\subsection{Vectorized Throughput}

The single-series result understates \kalbee{}'s intended use case at
scale. Table~\ref{tab:bench-vec} repeats the same task across 1{,}000
independent series --- e.g., 1{,}000 simultaneously tracked
objects --- comparing a naive per-series Python loop, \kalbee{}'s
purpose-built \texttt{VectorizedKalmanFilter} (which batches every filter
through the same \texttt{F @ x}, \texttt{F @ P @ F.T} NumPy call instead
of iterating), and simdkalman, whose entire value proposition is exactly
this scenario.

\begin{table}[t]
\centering
\caption{Vectorized throughput over 1{,}000 independent series.}
\label{tab:bench-vec}
\begin{tabular}{lr}
\toprule
\textbf{Approach} & \textbf{Total time (ms)} \\
\midrule
\kalbee{}, naive per-series loop        & 13{,}442 \\
simdkalman, batched                     & 728 \\
\kalbee{}, \texttt{VectorizedKalmanFilter} & \textbf{184} \\
\bottomrule
\end{tabular}
\end{table}

\kalbee{}'s vectorized filter is approximately $4\times$ faster than
simdkalman and $73\times$ faster than the naive loop over the same
1{,}000 series, while sharing the same \texttt{predict}/
\texttt{update}/\texttt{filter\_sequence} interface as every other filter
in the library --- a user does not adopt a separate vectorization library,
only a different \texttt{BaseFilter} subclass.

\subsection{Summary}

The evaluation supports a specific, non-universal claim: \kalbee{}
reproduces reference accuracy exactly, is competitive but not
best-in-class on bare single-filter overhead, and is the fastest option
evaluated for batched multi-series filtering. We consider this an honest
result more useful to a prospective adopter than an unqualified
``fastest'' claim would be.

\section{Discussion and Limitations}

\textbf{Typing coverage.} \texttt{py.typed} and a \texttt{mypy}
configuration are in place, and all newly added modules (the CLI, the
sensor-fusion cookbook, the scikit-learn adapter, the Jacobian utility)
type-check cleanly. Retrofitting strict typing across the $\sim$30
pre-existing filter and smoother modules --- most of whose gaps stem from
\texttt{Optional[np.ndarray]} attributes on \texttt{BaseFilter} not being
narrowed before use --- remains future work rather than a claim we make
here.

\textbf{No GPU/autodiff backend.} Unlike kalman-jax, \kalbee{} does not
currently offer a JAX or GPU-accelerated backend. This is a deliberate
scope decision rather than an oversight: a differentiable backend is a
parallel implementation surface, not an incremental feature, and is left
as future work pending demonstrated demand.

\textbf{Yaw observability.} As discussed in Section~\ref{sec:cookbook},
the shipped attitude EKF cannot observe yaw from accelerometer data alone;
this is a property of the sensing modality, not a filter defect, but it is
worth restating as a limitation of the \emph{recipe} as shipped, absent a
magnetometer or course-over-ground update.

\section{Conclusion and Future Work}

We presented \kalbee{}, a Python library that consolidates eighteen
Kalman-filter variants, four smoothers, multi-target tracking, parameter
learning, and two applied sensor-fusion recipes behind one consistent
interface, and evaluated it honestly against the closest prior art on a
reproducible benchmark. The results show \kalbee{} matches reference
accuracy, trades a modest amount of single-filter overhead for far greater
breadth, and substantially outperforms a comparable vectorization-focused
library at the batched-tracking workload that motivated it. Future work
includes a JAX-based autodiff/GPU backend for gradient-based noise-covariance
tuning and large-scale parallel filtering, a magnetometer-augmented
attitude filter for full observability, and completing the strict-typing
retrofit across the legacy filter modules.

\section*{Availability}
\kalbee{} is released under the Apache License 2.0 and is available on
PyPI (\texttt{pip install kalbee}) and at
\url{https://github.com/MinLee0210/kalbee}, with documentation at
\url{https://minlee0210.github.io/kalbee}.

\begingroup
\small
\begin{thebibliography}{00}

\bibitem{kalman1960} R. E. Kalman, ``A new approach to linear filtering and
prediction problems,'' \emph{Journal of Basic Engineering}, vol. 82, no. 1,
pp. 35--45, 1960.

\bibitem{julier2004unscented} S. J. Julier and J. K. Uhlmann, ``Unscented
filtering and nonlinear estimation,'' \emph{Proceedings of the IEEE},
vol. 92, no. 3, pp. 401--422, 2004.

\bibitem{arasaratnam2009cubature} I. Arasaratnam and S. Haykin, ``Cubature
Kalman filters,'' \emph{IEEE Transactions on Automatic Control}, vol. 54,
no. 6, pp. 1254--1269, 2009.

\bibitem{rauch1965maximum} H. E. Rauch, F. Tung, and C. T. Striebel,
``Maximum likelihood estimates of linear dynamic systems,'' \emph{AIAA
Journal}, vol. 3, no. 8, pp. 1445--1450, 1965.

\bibitem{bewley2016sort} A. Bewley, Z. Ge, L. Ott, F. Ramos, and B.
Upcroft, ``Simple online and realtime tracking,'' in \emph{Proc. IEEE
International Conference on Image Processing (ICIP)}, 2016, pp. 3464--3468.

\bibitem{sola2017quaternion} J. Sol\`{a}, ``Quaternion kinematics for the
error-state Kalman filter,'' \emph{arXiv preprint arXiv:1711.02508}, 2017.

\bibitem{filterpy} R. Labbe, ``FilterPy: Python Kalman filtering and
optimal estimation library,'' GitHub repository, 2014--present. [Online].
Available: \url{https://github.com/rlabbe/filterpy}

\bibitem{pykalman} pykalman contributors, ``pykalman: Kalman Filter,
Smoother, and EM Algorithm for Python,'' GitHub repository. [Online].
Available: \url{https://github.com/pykalman/pykalman}

\bibitem{simdkalman} O. Seiskari, ``simdkalman: Python Kalman filters
vectorized as Single Instruction, Multiple Data,'' GitHub repository.
[Online]. Available: \url{https://github.com/oseiskar/simdkalman}

\bibitem{stonesoup} Defence Science and Technology Laboratory (Dstl),
``Stone Soup: A software project to provide the target tracking community
with a framework for the development and testing of tracking algorithms,''
GitHub repository, 2017--present. [Online]. Available:
\url{https://github.com/dstl/Stone-Soup}

\bibitem{kalmanjax} AaltoML, ``kalman-jax: Approximate inference for
Markov Gaussian processes using iterated Kalman smoothing, in JAX,''
GitHub repository. [Online]. Available:
\url{https://github.com/AaltoML/kalman-jax}

\bibitem{kalbee} L. D. Minh, ``kalbee: A clean, modular Python
implementation of Kalman filters and estimation algorithms,'' v0.6.0,
2026. [Online]. Available: \url{https://github.com/MinLee0210/kalbee}

\end{thebibliography}
\endgroup

\end{document}
